\documentclass[10pt,a4paper]{article}
\usepackage[margin=1.8cm]{geometry}
\usepackage[T1]{fontenc}\usepackage{mathptmx}
\usepackage{amsmath,booktabs,tabularx,array,enumitem,xcolor}
\usepackage{xurl,needspace}\usepackage[hidelinks]{hyperref}
\usepackage{titlesec}
\definecolor{navy}{HTML}{1E2761}
\titleformat{\section}{\large\bfseries\color{navy}}{\thesection}{0.6em}{}
\titleformat{\subsection}{\normalsize\bfseries\color{navy}}{}{0em}{}
\titlespacing*{\section}{0pt}{1.1em}{0.5em}
\titlespacing*{\subsection}{0pt}{0.9em}{0.3em}
\setlength{\parindent}{0pt}\setlength{\parskip}{0.45em}
\newcolumntype{L}{>{\raggedright\arraybackslash}X}
\renewcommand{\arraystretch}{1.25}
\newcommand{\paper}[9]{%
\subsection*{#1}
\begin{tabularx}{\linewidth}{@{}>{\bfseries}p{3.4cm}L@{}}\toprule
Problem and method & #2\\
Contribution and data & #3\\
Baseline and result & #4\\
Limitations & #5\\
Colab feasibility & #6\\
Our next question & #7\\
Evidence & #8\\\bottomrule
\end{tabularx}}
\begin{document}
{\LARGE\bfseries\color{navy} HireMind Deep Learning Research Pack}\\[0.3em]
{\large Five papers and a connected experiment plan}\\[0.4em]
Rawan Khojah (S22107688) and Hams Aljohani (S20106833) \textbar{} CS4085 Deep Learning \textbar{} Effat University \textbar{} Item 1 preparation, October 2026

\section{One platform and one research component}
HireMind is our Senior project, an AI-assisted recruitment platform. This Deep Learning project studies one component it could use: a small model that receives CV text and a job description and predicts \emph{No Fit}, \emph{Potential Fit} or \emph{Good Fit}. After evaluation, its output could become one evidence source in the platform's scoring workflow. Practical challenges, explanations, storage and recruiter decisions remain separate platform functions.

The task is CV--job fit classification, not extracting individual skill names. Fit labels support the former; skill extraction would require different annotations and experiments.

\section{Reading notes: recruitment and language models}
Each note answers the assignment's paper questions. Results belong to the cited papers, not to HireMind. Feasibility judgments and next questions are our proposed adaptations.

\paper{Paper 1 \quad Gan, Zhang and Mori (2024)}
{Automate resume screening using LLM sentence classification, grading, summarization and agents.}
{Integrated screening study; 1,000 predominantly IT resumes with 78,668 classified sentences; separate grading labels largely generated by GPT-4.}
{Table 1: instruction-formatted LLaMA2-7B-chat reaches 87.73 F1, versus 86.67 for BERT Large.}
{Different tasks from our classifier. Table 10's speed comparison uses estimated human reading time. The grading data description gives 838 total versus 500 training plus 383 test resumes.}
{Full reproduction is unsuitable for our budget; we use the recruitment context, not its large-model training setup.}
{Can a compact classifier supply useful CV--job fit predictions at lower deployment cost?}
{Sections 4--6; Tables 1 and 10; PDF pages 8--16.}{}

\paper{Paper 2 \quad Devlin, Chang, Lee and Toutanova (2019)}
{Learn transferable bidirectional text representations using a Transformer, masked language modeling and next-sentence prediction.}
{One pretrained architecture can be fine-tuned for many tasks. Pretraining uses BooksCorpus and Wikipedia; evaluation includes GLUE and SQuAD.}
{Table 1: BERT Base scores 84.6\% on MNLI matched, compared with 82.1\% for OpenAI GPT.}
{General NLP results do not establish recruitment validity. Large pretraining cost and a 512-token input limit matter for long CV--job pairs.}
{Fine-tune an existing checkpoint; do not reproduce pretraining.}
{How much task quality can a compact BERT retain under a fixed input budget?}
{Sections 3--5, Table 1 and Appendix A; PDF pages 3--9 and 13--14.}{}

\needspace{14\baselineskip}\section{Compact models and distillation}
\paper{Paper 3 \quad Turc, Chang, Lee and Toutanova (2019)}
{Build effective compact language models through pretraining followed by task-specific teacher distillation.}
{Compares 24 model sizes and examines transfer-data size and domain. Tasks include GLUE and Amazon review classification.}
{Table 3: a 6-layer student's six-task development meta-score rises from 82.8 with pretraining plus fine-tuning to 84.4 with pretrained distillation.}
{The result is not for BERT-Tiny or recruitment. Benefits depend on task and transfer data; reported hardware speedups cannot be copied to our CPU.}
{Use a released compact checkpoint and one teacher; omit large-scale pretraining and the 24-model sweep.}
{Does distillation improve the same small classifier on limited CV--job pairs?}
{Sections 3--6; Tables 1, 3 and 4; PDF pages 3--9.}{}

\paper{Paper 4 \quad Hinton, Vinyals and Dean (2015)}
{Transfer knowledge from an expensive model or ensemble into a deployable student using softened class probabilities.}
{Temperature-based distillation, optionally combined with a hard-label loss. Experiments cover MNIST, speech and specialist image models.}
{On MNIST, the smaller hard-label model makes 146 errors; its distilled version makes 74; the larger teacher makes 67.}
{Not a BERT or recruitment study. Some exploratory missing-class experiments adjust biases using the test set; we must not copy that evaluation practice.}
{A small teacher--student adaptation is feasible; reproducing the large speech or image experiments is outside scope.}
{Do teacher probabilities improve generalization beyond the same student trained on hard labels?}
{Sections 2--4 and 8; PDF pages 2--5 and 8.}{}

\needspace{14\baselineskip}\section{Quantization and the combined research gap}
\paper{Paper 5 \quad Zafrir, Boudoukh, Izsak and Wasserblat (2019)}
{Reduce BERT deployment cost with 8-bit quantization-aware training, quantizing embedding and fully connected operations while keeping selected higher-precision operations.}
{Evaluates quantization on GLUE tasks and SQuAD; compares FP32, quantization-aware training and dynamic quantization.}
{MRPC F1: 90.00 FP32, 89.56 QAT and 87.88 dynamic quantization (development-set averages over five runs). About fourfold weight compression.}
{Hardware acceleration is conditional. Its QAT procedure differs from our post-training method; its fourfold result is not a HireMind measurement.}
{Use CPU dynamic Linear quantization plus weight-only embeddings (post-training); QAT stays outside the experiment.}
{Can the quantized student retain the distillation gain and reduce serialized model size?}
{Sections 2--4, Table 1; PDF pages 2--4.}{}

\subsection{What the five papers support together}
The recruitment paper motivates the application. BERT supplies the architecture foundation. Turc and Hinton motivate teacher-guided training. Q8BERT supports investigating precision reduction and highlights possible quality loss. Together they justify an experiment; they do not guarantee a positive result.

\subsection{Research question}
\emph{Can teacher-guided distillation improve a pretrained compact CV--job fit classifier, and can post-training quantization reduce its storage cost while retaining that improvement?}

\subsection{What changes in the platform if the experiment succeeds}
A backend endpoint could accept a CV--job text pair and return three class probabilities, a predicted class and a model version. The platform could display this result alongside its other evidence. It must not be presented as verified skills, a calibrated hiring probability or a final hiring decision.

\subsection{Two cautions}
Quantization reduces precision, not necessarily parameter count, token count or theoretical operation count. Distillation can add training cost even if the deployed student is inexpensive. These costs are reported separately; no unmeasured energy savings are claimed.

\section{Three candidate ideas and our selection}
All three ideas use CV--job fit labels, a pretrained BERT-Tiny baseline, job-grouped validation and the same final test policy. Numerical targets are exploratory.

\begin{tabularx}{\linewidth}{@{}p{2.7cm}LL@{}}\toprule
Idea & M0 $\rightarrow$ M1 $\rightarrow$ M2 & Targets and main difficulty\\\midrule
\textbf{A Teacher-guided training} (selected) & M0: hard-label fine-tuning. M1: same student plus teacher distillation. M2: INT8 Linear layers and 8-bit embedding weights in M1. & H1: +1 macro-F1 point. H2: at least 60\% smaller, retaining at least half the gain. Difficulty: teacher quality and extra training.\\
\textbf{B Class-weighted training} (backup) & M0: ordinary cross-entropy. M1: inverse-frequency weighted cross-entropy (training counts only). M2: quantized M1. & Same targets. Difficulty: weighting may hurt other classes; direct literature support is thinner.\\
\textbf{C Longer text context} (alternative) & M0: 128 tokens. M1: same architecture with 256 tokens. M2: quantized M1 at 256 tokens. & Same targets. Difficulty: longer input increases runtime and changes the input budget.\\\bottomrule
\end{tabularx}

\textbf{Why A.} It uses the selected papers most directly and isolates a training change while keeping the student architecture and input length fixed. It could produce a small deployable component for the platform. B is cheaper if the teacher is ineffective; C risks making the model slower before compression.

\subsection{Model roles}
\begin{tabularx}{\linewidth}{@{}>{\bfseries}p{2.6cm}L@{}}\toprule
Teacher & \texttt{google-bert/bert-base-uncased} with a three-class head, fine-tuned only on the training portion; checkpoint chosen on validation; training logits cached.\\
M0 & \texttt{google/bert\_uncased\_L-2\_H-128\_A-2} (about 4.4M parameters) with a three-class head, trained on hard labels.\\
M1 & Same pretrained student initialization and head seed as M0, trained with $0.5\cdot\mathrm{CE} + 0.5\cdot T^2\cdot \mathrm{KL}(p_{\text{teacher}}\,\|\,p_{\text{student}})$, both softened at temperature $T$.\\
M2 & CPU copy of each M1: dynamic INT8 Linear modules plus 8-bit weight-only Embedding modules; LayerNorm and unsupported operations stay FP32; no additional training.\\\bottomrule
\end{tabularx}

M1 does not continue from trained M0: both start from matched initializations with equal update budgets. The teacher supports M1; it does not replace the M0 baseline.

\section{Dataset audit and evaluation boundary}
Dataset: \texttt{cnamuangtoun/resume-job-description-fit}, revision\\\texttt{08978e21714984bb417547d2c0f9b477f5298163}, approved by the instructor. File-structure checks (13 October 2026) are not model evaluation.

\begin{tabularx}{\linewidth}{@{}>{\bfseries}p{3.4cm}L@{}}\toprule
Schema & \texttt{resume\_text}, \texttt{job\_description\_text}, \texttt{label}\\
Original split & 6,241 training rows; 1,759 test rows; 8,000 total\\
Labels & No Fit 4,000; Potential Fit 2,000; Good Fit 2,000\\
Unique raw text & 643 CVs and 351 job descriptions across both files\\
Repeated pairs & 7,993 unique raw-text pairs; seven excess pair occurrences\\
Conflicting labels & Six identical pairs carry conflicting labels (twelve rows)\\
Split overlap & 476 identical CV texts occur in both files; no identical job descriptions overlap\\
Empty fields & None in the three required columns\\
Documentation & No dataset card or licence in this revision; label origin unverified\\\bottomrule
\end{tabularx}

\textbf{Cleaning and splitting plan.} Remove the twelve conflicting rows and collapse the remaining duplicate, keep the official test jobs, and reserve about 20\% of training-job groups for validation (seed 42) so each job stays in one split. Check near duplicates before freezing the split; never choose the split using model performance.

\textbf{What the test can and cannot show.} It evaluates new job descriptions with potentially familiar CVs. It does not measure generalization to unseen candidates, real hiring outcomes or fairness across demographic groups. Raw CV text stays out of the repository.

\section{Experiment and hypotheses}
\textbf{Fair training.} Teacher, M0 and M1 share labels, tokenizer, 256-token maximum length, longest-first pair truncation and preprocessing. AdamW, effective batch 16, weight decay 0.01, at most three epochs; checkpoints selected by validation macro-F1. Equal search budget: M0 learning rates $\{2,5,10,30\}\times10^{-5}$; M1 learning rates $\{5,10\}\times10^{-5}$ $\times$ temperatures $\{2,4\}$ with loss weight 0.5; teacher $\{2,5\}\times10^{-5}$. Search at seed 42, then rerun the chosen settings at seeds 43 and 44.

\textbf{Hypotheses.} With $F_0,F_1,F_2$ the mean test macro-F1 over paired seeds and $S_1,S_2$ serialized sizes:
H1: $F_1-F_0\ge 1.0$ point. H2: $S_2\le0.40\,S_1$, $F_1-F_2\le1.0$ point and $F_2-F_0\ge0.5(F_1-F_0)$ with $F_1>F_0$. If M1 does not beat M0, retention cannot be confirmed and compression is reported separately.

\begin{tabularx}{\linewidth}{@{}>{\bfseries}p{3.4cm}L@{}}\toprule
Quality & Macro-F1 primary; accuracy, per-class precision/recall/F1 and confusion matrix secondary. Choices frozen before the test.\\
Size & Same serialization for all models, excluding optimizer state; bytes and MiB.\\
CPU latency & Same CPU, one thread, batch 1; 30 warm-up and 200 timed calls; median and p95, plus a tokenization-inclusive measurement.\\
Memory & Fresh process per model; peak resident memory; GPU training peak where available.\\
MACs & Same input shape and counting convention; quantization can leave operation counts unchanged.\\
Uncertainty and cost & Paired-seed mean and standard deviation; teacher, student and logit-cache time recorded separately.\\\bottomrule
\end{tabularx}

\needspace{4\baselineskip}\textbf{Risks and fallbacks.} Weak teacher: report negative transfer. Unsupported embedding quantization: fall back to Linear-only quantization and revise H2 before the test. Quality loss: report it as measured.

\section{Sources}
\begin{enumerate}[leftmargin=1.6em,itemsep=0.2em,label={[\arabic*]}]
\item Chengguang Gan, Qinghao Zhang and Tatsunori Mori. 2024. Application of LLM Agents in Recruitment: A Novel Framework for Resume Screening. arXiv:2401.08315v2. \url{https://arxiv.org/abs/2401.08315v2}
\item Jacob Devlin, Ming-Wei Chang, Kenton Lee and Kristina Toutanova. 2019. BERT: Pre-training of Deep Bidirectional Transformers for Language Understanding. NAACL-HLT, pp.~4171--4186. \url{https://aclanthology.org/N19-1423/}
\item Iulia Turc, Ming-Wei Chang, Kenton Lee and Kristina Toutanova. 2019. Well-Read Students Learn Better: On the Importance of Pre-training Compact Models. arXiv:1908.08962v2. \url{https://arxiv.org/abs/1908.08962v2}
\item Geoffrey Hinton, Oriol Vinyals and Jeff Dean. 2015. Distilling the Knowledge in a Neural Network. arXiv:1503.02531.\\ \url{https://arxiv.org/abs/1503.02531}
\item Ofir Zafrir, Guy Boudoukh, Peter Izsak and Moshe Wasserblat. 2019. Q8BERT: Quantized 8Bit BERT. EMC2 workshop. DOI: 10.1109/EMC2-NIPS53020.2019.00016. \url{https://arxiv.org/abs/1910.06188v2}
\end{enumerate}
\textbf{Resources.} Dataset: \url{https://huggingface.co/datasets/cnamuangtoun/resume-job-description-fit} \textbullet{} Student: \url{https://huggingface.co/google/bert_uncased_L-2_H-128_A-2} \textbullet{} Teacher: \url{https://huggingface.co/google-bert/bert-base-uncased} \textbullet{} PyTorch dynamic quantization: \url{https://docs.pytorch.org/docs/main/generated/torch.ao.quantization.quantize_dynamic.html}
\end{document}
